\section{Aturan Resolusi dalam Kecerdasan Buatan}

Seiring bergulirnya revolusi komputer abad modern, ahli komputasi mengidamkan suatu metode mesin mandiri yang sanggup mengeksekusi inferensi jutaan fakta tanpa harus menghafal kesebelas varian Modus Ponens/Tollens secara spesifik \cite{rosen2019,iep_natded}. John Alan Robinson (1965) menjawab kebutuhan ini melalui algoritma penalaran yang menyatukan deduksi: \textbf{Resolusi}.

\subsection{Definisi dan Mekanika Resolusi}

Resolusi adalah teknik inferensi silang pada dua buah klausa Disjungsi ($OR$) yang memuat satu pasangan \textit{Proposisi Saling Bertentangan (Literal dan Negasinya)}.

\begin{teorema}[Prinsip Resolusi]
Diberikan dua klausa disjungsi: $(p \lor q)$ dan $(\neg p \lor r)$. Berdasarkan Hukum Non-Kontradiksi, $p$ dan $\neg p$ tidak dapat benar bersamaan. Jika $p$ benar, maka $\neg p$ salah sehingga $r$ harus benar (dari klausa kedua). Jika $p$ salah, maka $q$ harus benar (dari klausa pertama). Konklusinya: $(q \lor r)$---gabungan literal sisa yang disebut \textbf{resolvent}.
\end{teorema}

\textbf{Anatomi Blok Argumen Resolusi:}
\[
\begin{array}{cl}
p \vee q & \text{(Klausa OR Pertama)} \\
\neg p \vee r & \text{(Klausa OR Kedua memuat penangkal } p) \\
\hline
\therefore q \vee r & \text{(Resolvent / Kesimpulan Hasil Silang)}
\end{array}
\]

\begin{contoh}
Misalkan dua fakta logis dikirimkan ke dalam mesin Kecerdasan Buatan (AI):
\begin{itemize}
    \item \textbf{Data 1 ($p \lor q$)}: "Andi sedang mudik menginjak tanah Surabaya ATAU Andi asyik bertugas piket di ibukota Jakarta."
    \item \textbf{Data 2 ($\neg p \lor r$)}: "Andi jelas Tidak Sedang Berada di tanah mudikan Surabaya ATAU mesin mobil operasionalnya kini tengah berjalan panik mengular di tol malam."
\end{itemize}
Sistem memindai pasangan komplementer $p$ dan $\neg p$, mengeliminasinya, dan menggabungkan sisa literal: \\
\textbf{Konklusi ($\therefore q \lor r$)}: ``Andi bertugas piket di Jakarta ($q$) ATAU mobilnya berjalan di tol ($r$).''
\end{contoh}

Dalam arsitektur *Artificial Intelligence* masa kini, khususnya sistem pakar bahasa pemrograman \textbf{PROLOG (Programming in Logic)}, hampir 100\% motor inferensi di balik layarnya berjalan melumat basis data (*Knowledge Base*) memakai repetisi mesin potong \textit{Resolusi} berantai ini \cite{wiki_first_order_logic,libretexts_proofs}. Sangat elegan dan mudah dikodekan dalam loop biner.
